\documentclass[10pt]{article}
\usepackage[letterpaper,margin=.38in]{geometry}
\usepackage{amsmath,amssymb,multicol,graphicx}
\pdfinfoomitdate=1\pdftrailerid{}
\pagestyle{empty}\setlength{\parindent}{0pt}\setlength{\parskip}{2pt}
\setlength{\columnsep}{13pt}\setlength{\jot}{0pt}
\setlength{\abovedisplayskip}{3pt}\setlength{\belowdisplayskip}{3pt}
\let\bigsum\sum\renewcommand{\sum}{\bigsum\nolimits}
\newcommand{\fl}[1]{\left\lfloor#1\right\rfloor}
\newcommand{\rg}[2]{\{#1,\ldots,#2\}}
\begin{document}
\begin{center}
{\large\bfseries A thirty-two-variable consistency equation}\\[-1pt]
{\small Positive integer unknowns; exact degree 6; a 142-operation evaluator.}
\end{center}\vspace{-8pt}
\begin{multicols}{2}
\fontsize{8.6}{9.1}\selectfont
% BEGIN VERIFIED DISPLAY
The only unknowns are the thirty-two separately named positive-integer variables
\[
\begin{gathered}J,F,\alpha,f,g,T,X,Y,\xi,\eta,\zeta,\upsilon,\gamma,\psi,\varphi,\epsilon,\\
c,\ell,W,m,n,Z,D,e,z,\delta,\\
\kappa,\lambda,\tau,\chi,\theta,Q\in\mathbb Z_{>0}.\end{gathered}
\]
All other symbols below are polynomial abbreviations:
\[
\begin{gathered}
b=\beta J,\quad q=1+b,\quad C=q-F-\alpha,\quad U=Y\kappa,\\
a=\lambda+Y,\quad s=q\kappa,\quad x=\lambda q,\\
\sigma=\psi^2,\quad\nu=\upsilon^2,\\
E=\lambda,\quad k=\eta+\zeta,\\
H=3+4a,\quad\Delta=\delta,\quad L=X+ac,\\
M=sb,\\
R=[\alpha+(q-F)b](\kappa-1)+J(\mu+\omega q),\\
V=e-z,\\
\vartheta=V^2,\\
I=\theta,\\
\mathcal{P}=Z^2-\sigma(1+\tau),\\
\mathcal{D}=D^2-\chi(1+\Delta n),\\
\mathcal{Q}=\nu-Q[\nu-\vartheta]-1,\\
\mathcal{R}=I^2-Q,\\
\mathcal{B}=\psi-M,\quad \mathcal{C}=\varphi-U,\quad \mathcal{G}=\epsilon-R,\\
\mathcal{H}=c-\eta-Yk,\\
\mathcal{L}=\ell-EYk,\quad \mathcal{W}=W-[C(X+Kq)-Fq],\\
\mathcal{Z}=m-f^2,\quad \mathcal{N}=n-c^2,\\
\mathcal{I}=Z-[\psi(\ell+k+\xi-kW)+Yb-Ws],\\
\mathcal{J}=D-[\varphi L+H(\varphi\gamma+x+\epsilon+1-k)],\\
\mathcal{K}=V-[m(Tc-\epsilon)-c],\quad\mathcal{O}=\delta-H-a^2,\\
\mathfrak{J}=\kappa-q^2,\quad\mathcal{F}=\lambda-XY,\\
\mathcal{S}=\tau-\ell(\ell+k),\\
\mathcal{U}=\chi-\varphi^2,\quad\mathcal{V}=\theta-ng,\\
\mathfrak{B}=Q-\Delta(m-1).
\end{gathered}
\]
The single Diophantine equation is
\[
\begin{split}
\mathcal{P}^2+\mathcal{D}^2+\mathcal{Q}^2+\mathcal{R}^2+\mathcal{B}^2+\mathcal{C}^2+\mathcal{G}^2 &\\
+\mathcal{H}^2+\mathcal{L}^2+\mathcal{W}^2+\mathcal{Z}^2+\mathcal{N}^2 &\\
+\mathcal{I}^2+\mathcal{J}^2+\mathcal{K}^2+\mathcal{O}^2+\mathfrak{J}^2+\mathcal{F}^2 &\\
+\mathcal{S}^2+\mathcal{U}^2+\mathcal{V}^2+\mathfrak{B}^2 &= 0.
\end{split}
\]
% END VERIFIED DISPLAY
\textbf{The four fixed integers.}
Every index, set, sum, product and floor below defines a constant, not
another unknown. Put
\[
\begin{gathered}
\rho=2^{47},\quad B=\rho^{2011},\quad\beta=B-1,\\
G_m(x)=\sum_{j=0}^{m-1}x^j,\quad
\Lambda=\rho^{1005}.
\end{gathered}
\]
\begin{align*}
\mu={}&\beta-G_{951}(\rho)\\
&-\rho^{960}G_{26}(\rho)-\rho^{1004},\\
\omega={}&\beta+8G_{917}(\rho)-4+(2^{45}-4)\Lambda\\
&+\rho^{960}G_9(\rho)+\rho^{987}G_8(\rho)
  +\rho^{1004},\\
K={}&\rho^{18}+\rho^{1004}+\rho^{2020}\\
&+2\Lambda\Phi(1/\rho),
\end{align*}
where
\[
\begin{gathered}
\chi_i=1+2^{10}(1+w_i+2^8n_i\\
\qquad+2^{17}p_i+2^{25}s_i)\quad(0\le i\le941),\\
\pi(0)=0,\quad\pi(i)=1+\sum_{j=1}^{941}
\fl{\frac{2^{44}-1+\chi_j-\chi_i}{2^{44}}}\quad(1\le i\le941).
\end{gathered}
\]
\[
\begin{split}
\Phi(x)={}&\sum_{i=0}^{941}\chi_i x^{\pi(i)}\\
&+2^{10}x^{969}(x^9+2^{17})G_8(2x)\\
&+2^{18}x^{942}(x^{18}+2^{17})G_9(2x).
\end{split}
\]
The rational intermediate in K cancels exactly; all four results are positive integers.

\textbf{The 942 indexed fixed tiles.}
Use the actions $u_j,v_j,d_j$ decoded below. For $d=0,1$, let
\[
\begin{gathered}
h(q,b)=2q+b+2,\quad
\mathcal S_d=\{u_j:d_j=d,\ 0\le j<430\},\\
r_{d,q}=205-199d+|\mathcal S_d\cap\rg0{q-1}|.
\end{gathered}
\]
An interval ending at $-1$ is empty. The table defines $t_i$ for
$0\le i\le941$; all index ranges are fixed:
\[
\setlength{\arraycolsep}{1.4pt}\begin{array}{c|l}
\text{index }i&\text{tile }t_i=(w_i,n_i,p_i,s_i)\\\hline
0&(2,435,3,1)\\
j+1\ (j=2q+b<430)&((1-d_j)r_{0,u_j},h(q,b),d_jr_{1,u_j},v_j)\\
433-2a+b\ (a,b=0,1)&(0,432a+b,0,432a+b)\\
\substack{423+2r_{d,q}+b\\q\in\mathcal S_d,\ d,b=0,1}&(dr_{d,q},b,(1-d)r_{d,q},h(q,b))\\
923+2a+b\ (a,b=0,1)&(0,432a+b,0,438)\\
927&(2,434,2,0)\\
928&(3,436,4,4)\\
929&(4,436,4,0)\\
496+t\ (t=434,435,436)&(1,438,1,t)\\
933+3a+b\ (a,b=0,1,2)&(a^2,437,b,437)
\end{array}
\]
\textbf{A closed decoder for the 430 actions.}
Put $D_b(A,k)=\fl{A/b^k}\bmod b$, where
$a\bmod b=a-b\fl{a/b}$. For $0\le q<215$, $b=0,1$, $j=2q+b$:
\[
\begin{gathered}
\mathcal N=\mathcal A\bmod2^{636},\quad
\mathcal T=\fl{\mathcal A/2^{636}},\\
\nu_j=D_2(\mathcal N,j),\quad \kappa_j=(\sigma_j+2)^{1-\nu_j}\quad(0\le j<430),\\
\sigma_j=\sum_{k=0}^{j-1}\nu_k,\quad
\Pi_j=\prod_{k=0}^{j-1}\kappa_k\quad(0\le j\le430),\\
m_q=D_2(\mathcal N,430+q),\quad \tau_q=\sum_{k=0}^{q-1}m_k,\\
F_q=14+m_q(D_{16}(\fl{\mathcal T/\Pi_{430}},\tau_q)-14),\\
u_j=(\sigma_{j+1}-(\fl{\mathcal T/\Pi_j}\bmod\kappa_j))\bmod216,\\
v_j=D_2(F_q,2b),\quad d_j=D_2(F_q,2b+1).
\end{gathered}
\]
Empty sums/products are $0/1$. Here $(v_0,d_0,u_0)=(1,1,1)$; the constructor tile
uses this actual first target, including its state renaming.

\textbf{The integer $\mathcal A$ in base 10.}
Its complete decimal expansion is the concatenation of these lines:\par
\begingroup\fontsize{7.7}{8.2}\selectfont\ttfamily
12889260410359561746736540743515845134244260640012702033212621\par
95398074180784058571518488282119861023093355389406939414042892\par
97455773992193784795705001618777025966369137954859764532832722\par
75411843410809038713502633204617333246733616940964578272360901\par
56186216508504953649192752154519268024509952082950238865468307\par
22999632482769253980274091217460707829023970817244790340704449\par
09085581271192071656170732004937337082171141082486384269072884\par
72778087419663417942892520175630649792329265834490253035577533\par
96449689808375788457488971017043689875433162972390177936123032\par
89025657757489526342497438772045714099661532481123050235898973\par
27443609903800541199351405985754278307227777584454508422717998\par
334113564409335346074892680328437\par
\endgroup
\normalfont\fontsize{7.7}{8.2}\selectfont
The fixed integer $\mathcal A$ above is written in base 10. Lean proves that the formalized equation has no positive solution exactly
when the chosen arithmetized ZFC consistency statement holds, by a derivation
in the specified intuitionistic PRA. The literal endpoint was independently
replayed on a fixed source snapshot; some transitive dependencies remain
inherited. The exact degree is certified by integer-polynomial calculations;
the 142-operation count is for the unshifted circuit. Neither count is a
minimality claim.
\end{multicols}
\clearpage
\begin{center}{\large\bfseries Exact 142-operation circuit}\end{center}
Each numbered line defines one gate over the integers. The inputs are the
thirty-two unknowns and the four fixed integers above. The ASCII names
\texttt{alpha, xi, eta, zeta, upsilon, gamma, vartheta, psi, varphi, epsilon, beta, mu, omega, ell, Delta, delta, kappa, lambda, sigma, tau, chi, theta, nu}
refer to their displayed counterparts; \texttt{g} is the unknown $g$. The last
gate is zero exactly when the displayed equation holds. The circuit inputs
\texttt{cNamed, ellNamed} are $c,\ell$; gate \texttt{a} is the displayed abbreviation; input \texttt{WNamed} is the
displayed $W$, while gate \texttt{WComputed} is $C(X+Kq)-Fq$. The inputs
\texttt{m, n, Z, D, e, z, delta} are the displayed
$m,n,Z,D,e,z,\delta$; \texttt{kappa, lambda, tau, chi, theta, Q}
are $\kappa,\lambda,\tau,\chi,\theta,Q$. Gates \texttt{q3, g55, g39, y2, g80} compute $s,x,\sigma,\nu,\vartheta$. Gates \texttt{E, I} compute
$XY,ng$; the display uses $E=\lambda$, $I=\theta$. Intermediate gate
names below are local abbreviations, not additional unknowns.
\begin{multicols}{3}
\raggedright\ttfamily\fontsize{7.2}{8.4}\selectfont
1. b = J * beta\\
2. q = 1 + b\\
3. q2 = q * q\\
4. q3 = q * kappa\\
5. U = Y * kappa\\
6. qF = q - F\\
7. C = qF - alpha\\
8. E = X * Y\\
9. k = eta + zeta\\
10. kY = Y * k\\
11. c = eta + kY\\
12. c2 = cNamed * cNamed\\
13. a = lambda + Y\\
14. g15 = 4 * a\\
15. H = 3 + g15\\
16. g17 = a * a\\
17. DeltaComputed = H + g17\\
18. g19 = a * cNamed\\
19. L = X + g19\\
20. ell = lambda * kY\\
21. M = q3 * b\\
22. g25 = K * q\\
23. g26 = X + g25\\
24. g27 = C * g26\\
25. g28 = F * q\\
26. WComputed = g27 - g28\\
27. g30 = ellNamed + k\\
28. g31 = g30 + xi\\
29. root\_skew = k * WNamed\\
30. root\_inner = g31 - root\_skew\\
31. root\_scaled = psi * root\_inner\\
32. g33 = Y * b\\
33. root\_btq3 = WNamed * q3\\
34. root\_sum = root\_scaled + g33\\
35. ZComputed = root\_sum - root\_btq3\\
36. g38 = Z * Z\\
37. g39 = psi * psi\\
38. g40 = ellNamed * g30\\
39. g41 = 1 + tau\\
40. g42 = g39 * g41\\
41. A0 = g38 - g42\\
42. g44 = qF * b\\
43. rankseed = alpha + g44\\
44. g46 = kappa - 1\\
45. g47 = rankseed * g46\\
46. g48 = omega * q\\
47. g49 = mu + g48\\
48. g50 = J * g49\\
49. R = g47 + g50\\
50. g52 = k - epsilon\\
51. r0 = g52 - 1\\
52. g54 = varphi * gamma\\
53. g55 = lambda * q\\
54. g56 = g54 + g55\\
55. shift = g56 - r0\\
56. g58 = varphi * L\\
57. g59 = H * shift\\
58. DstarComputed = g58 + g59\\
59. g61 = D * D\\
60. g62 = varphi * varphi\\
61. g63 = delta * n\\
62. g64 = 1 + g63\\
63. g65 = chi * g64\\
64. A1 = g61 - g65\\
65. f2 = f * f\\
66. rho = m - 1\\
67. QComputed = delta * rho\\
68. I = n * g\\
69. g71 = theta * theta\\
70. g72 = g71 - Q\\
71. g74 = T * cNamed\\
72. g75 = g74 - epsilon\\
73. g76 = m * g75\\
74. VComputed = g76 - cNamed\\
75. V = e - z\\
76. y2 = upsilon * upsilon\\
77. g80 = V * V\\
78. g81 = y2 - g80\\
79. g82 = Q * g81\\
80. Fa = y2 - g82\\
81. auxResidual = Fa - 1\\
82. firstSquare = A0 * A0\\
83. mainSquare = A1 * A1\\
84. auxSquare = auxResidual * auxResidual\\
85. strongSquare = g72 * g72\\
86. firstMain = firstSquare + mainSquare\\
87. firstMainAux = firstMain + auxSquare\\
88. sumSquares = firstMainAux + strongSquare\\
89. scaleDiff = psi - M\\
90. scaleCheck = scaleDiff * scaleDiff\\
91. final = sumSquares + scaleCheck\\
92. mainDiff = varphi - U\\
93. mainCheck = mainDiff * mainDiff\\
94. result16 = final + mainCheck\\
95. rankDiff = epsilon - R\\
96. rankCheck = rankDiff * rankDiff\\
97. result17 = result16 + rankCheck\\
98. cDifference = cNamed - c\\
99. cConstraint = cDifference * cDifference\\
100. ellDifference = ellNamed - ell\\
101. ellConstraint = ellDifference * ellDifference\\
102. wDifference = WNamed - WComputed\\
103. wConstraint = wDifference * wDifference\\
104. addC = result17 + cConstraint\\
105. addEll = addC + ellConstraint\\
106. finalWithNamedValues = addEll + wConstraint\\
107. fSquareDifference = m - f2\\
108. fSquareConstraint = fSquareDifference * fSquareDifference\\
109. cSquareDifference = n - c2\\
110. cSquareConstraint = cSquareDifference * cSquareDifference\\
111. addFSquare = finalWithNamedValues + fSquareConstraint\\
112. addCSquare = addFSquare + cSquareConstraint\\
113. zDifference = Z - ZComputed\\
114. zConstraint = zDifference * zDifference\\
115. dDifference = D - DstarComputed\\
116. dConstraint = dDifference * dDifference\\
117. vDifference = V - VComputed\\
118. vConstraint = vDifference * vDifference\\
119. deltaDifference = delta - DeltaComputed\\
120. deltaConstraint = deltaDifference * deltaDifference\\
121. addZConstraint = addCSquare + zConstraint\\
122. addDConstraint = addZConstraint + dConstraint\\
123. addVConstraint = addDConstraint + vConstraint\\
124. addDELTAConstraint = addVConstraint + deltaConstraint\\
125. qSquareDifference = kappa - q2\\
126. qSquareConstraint = qSquareDifference * qSquareDifference\\
127. product2Difference = lambda - E\\
128. product2Constraint = product2Difference * product2Difference\\
129. ellProductDifference = tau - g40\\
130. ellProductConstraint = ellProductDifference * ellProductDifference\\
131. varphiSquareDifference = chi - g62\\
132. varphiSquareConstraint = varphiSquareDifference * varphiSquareDifference\\
133. scaledIndexDifference = theta - I\\
134. scaledIndexConstraint = scaledIndexDifference * scaledIndexDifference\\
135. qSignedDifference = Q - QComputed\\
136. qSignedConstraint = qSignedDifference * qSignedDifference\\
137. add\_qSquare = addDELTAConstraint + qSquareConstraint\\
138. add\_product2 = add\_qSquare + product2Constraint\\
139. add\_ellProduct = add\_product2 + ellProductConstraint\\
140. add\_varphiSquare = add\_ellProduct + varphiSquareConstraint\\
141. add\_scaledIndex = add\_varphiSquare + scaledIndexConstraint\\
142. add\_qSigned = add\_scaledIndex + qSignedConstraint\\
\end{multicols}
\end{document}
